% ------------------------------------------------------------------ 7
\section{Experimental setup and overall results}\label{sec:experiments}

\paragraph{Machine and build.} \val{env.cpu}, \val{env.cores} logical cores, \val{env.mem}\,GB RAM, \val{env.os}; \val{env.gxx},
flags \texttt{-std=c++17 -O2 -march=native}. The main table ran six solver processes in parallel
(\texttt{scripts/run\_bench.sh}); the heap comparison of \cref{sec:ablation} ran one process at a time and keeps the best of three
runs, because it compares times. The exact solvers have a 300\,s limit per instance.

\paragraph{Instances.} SteinLib sets B (18 sparse random graphs, \(n\le100\)), C (20 graphs with \(n=500\), up to \(12{,}500\) edges and
\(k\) up to 250), P6E (15 Euclidean-weight graphs, \(n=100\)) and I080 (100 graphs with \(n=80\) and Duin's incidence weights, which
are designed to make reduction tests ineffective), all with published optima. The generated sets are the adversarial instances of \cref{sec:counter}, 27 grids
(\(10\times10\), \(k=4,6,\dots,20\), three each) and six sparse random graphs with \(n\) from 100 to 3200 and \(k=10\). In total
\val{inst.total} instances and \val{rows.total} solver runs.

\paragraph{Correctness checks.} Every tree any solver returns passes through the validator in \texttt{include/graph.hpp},
which is separate from all solvers: edge ids valid and distinct, no cycle, connected, all terminals spanned, and the cost recomputed from
the input. No run failed this check. \texttt{scripts/make\_tables.py} also refuses to build the report if any greedy cost is below the
known optimum, or if an exact solver that claims optimality differs from it. All \val{exact.checked} exact SteinLib results agree with the
published optima, and on the generated sets the DP and the branch-and-bound, two unrelated algorithms,
agree with each other. Finally, \texttt{scripts/fuzz.py} runs all five solvers on 300 random graphs with up to 14 vertices,
zero-weight edges, ties and \(k\) from 1 to \(n\); the two exact methods agree on all of them and no greedy tree is invalid or below
the optimum.

\begin{table}[!htbp]\centering\small
\begin{tabular}{@{}llrrrr@{}}\toprule
Solver & Paradigm & runs & optimal & mean gap (\%) & total time (s)\\\midrule
G1 & greedy & 199 & 18 & 14.195 & 0.7\\
G2 & greedy & 199 & 63 & 8.927 & 35.7\\
G3 & greedy & 199 & 60 & 4.646 & 3.3\\
DP & dynamic programming & 173 & 173 & 0.000 & 4617.3\\
B\&B & branch-and-bound & 199 & 198 & 0.018 & 379.3\\
\bottomrule\end{tabular}
\caption{All instances with a known optimum (199). \emph{runs}: runs that returned a tree (the DP skips \(k>20\)). For the exact methods a non-zero mean gap comes only from instances stopped at the time limit.}\label{tab:overview}
\end{table}


\subsection{Fibonacci heap against binary heap}\label{sec:ablation}
\begin{table}[H]\centering\small
\begin{tabular}{@{}llrrrrr@{}}\toprule
Solver & graphs & instances & Fibonacci (s) & binary (s) & Fib.\,/\,bin. & Fib. faster\\\midrule
G1 & sparse, \(m/n<10\) & 77 & 0.024 & 0.010 & 2.33 & 0/77\\
G1 & dense, \(m/n\ge10\) & 7 & 0.003 & 0.002 & 1.51 & 0/7\\
G2 & sparse, \(m/n<10\) & 77 & 0.146 & 0.056 & 2.60 & 0/77\\
G2 & dense, \(m/n\ge10\) & 7 & 0.010 & 0.006 & 1.58 & 0/7\\
G3 & sparse, \(m/n<10\) & 77 & 0.027 & 0.013 & 2.05 & 0/77\\
G3 & dense, \(m/n\ge10\) & 7 & 0.003 & 0.002 & 1.50 & 0/7\\
DP & sparse, \(m/n<10\) & 59 & 3.573 & 2.022 & 1.77 & 0/59\\
DP & dense, \(m/n\ge10\) & 7 & 0.072 & 0.055 & 1.31 & 0/7\\
\bottomrule\end{tabular}
\caption{Total time with each heap (best of three serial runs per instance). Sparse: SteinLib B, I080 (\(k\le16\) for the DP), grids and random graphs; dense: SteinLib c11, c12, c16, c17 (\(m/n=5\) to \(25\)) with c01, c02, c06, c07 as sparse C graphs. The DP returns the same optimum with both heaps; the greedy strategies can break ties between equal keys differently.}\label{tab:ablation}
\end{table}

\noindent The binary heap is faster on every instance, by a factor of \val{abl.DP.sparse.ratio} for the DP on sparse graphs and
\val{abl.G2.sparse.ratio} for G2. This does not contradict \cref{thm:sequence}: the theorem bounds the number of pointer
operations, not the time each one takes. The binary heap keeps its keys in one array that stays in cache, while every Fibonacci
node is a separate allocation with four pointers. On these graphs the asymptotic advantage has little to work with, because a
Dijkstra run on a graph with \(m/n<10\) does fewer decrease-keys than extract-mins (\cref{tab:heapcounts}), and each binary-heap
decrease-key moves a key only a few levels.

The gap does narrow where the theory says it should. On the dense graphs the DP ratio falls to \val{abl.DP.dense.ratio}, and
on \texttt{c17} (\(m/n=25\)) the run makes \val{abl.dp.c17.decr} decrease-keys against \val{abl.dp.c17.extr} extract-mins and
takes \val{abl.dp.c17.fib}\,ms with the Fibonacci heap against \val{abl.dp.c17.bin}\,ms with the binary heap. We report the
Fibonacci heap as the structure that gives the proved bounds, not as the faster choice on SteinLib-sized inputs.


\section{Limitations}\label{sec:limits}
\begin{itemize}
\item The DP is limited to \(k\le20\) by time and memory (\cref{thm:dpcost}). A meet-in-the-middle or the Nederlof
inclusion--exclusion variant would trade time for polynomial space, but would no longer use the heap.
\item The branch-and-bound has no LP solver, so its bound is the dual ascent value, which is weaker than the LP optimum on some dense
graphs. The instances left open in \cref{tab:bnb} are those where the root gap stays large; LP-based branch-and-cut solvers such as
SCIP-Jack are the standard tool for them.
\item The search for adversarial instances only covers graphs with at most 9 vertices. It shows that each greedy strategy can be
suboptimal, not how bad it can be: the ratios found (at most \val{adv.g1.ratio}) are far from the proven bounds.
\item Times are wall-clock times on one machine, measured inside each process from after the input is read until the solver returns its tree;
runs in the main table shared the machine with five others.
\end{itemize}
